\documentclass[10pt,landscape,twocolumn]{article}

% Page setup
\usepackage[landscape,top=1in,bottom=0.5in,left=0.5in,right=0.5in]{geometry}
\usepackage[utf8]{inputenc}
\usepackage[english]{babel}

% Math packages for formulas
\usepackage{amsmath}
\usepackage{amssymb}
\usepackage{amsthm}
\usepackage{mathtools}
\usepackage{amsfonts}
\usepackage{mathrsfs}

% Code listings
\usepackage{listings}
\usepackage{xcolor}

% Headers and footers
\usepackage{fancyhdr}
\usepackage{lastpage}

% Index/Table of contents
\usepackage{tocloft}
\usepackage{hyperref}

% Graphics
\usepackage{graphicx}

% Multiple columns for title page
\usepackage{multicol}

\usepackage{tabularx}     
\usepackage{ragged2e}
\usepackage{xltabular} 
\usepackage{multirow}   

% Listings configuration for C++
\lstset{
    language=C++,
    basicstyle=\ttfamily\footnotesize,
    keywordstyle=\color{blue}\bfseries,
    commentstyle=\color{green!50!black}\itshape,
    stringstyle=\color{red},
    numbers=left,
    numberstyle=\tiny\color{gray},
    stepnumber=1,
    numbersep=5pt,
    backgroundcolor=\color{white},
    showspaces=false,
    showstringspaces=false,
    showtabs=false,
    frame=single,
    tabsize=2,
    captionpos=b,
    breaklines=true,
    breakatwhitespace=false,
    escapeinside={\%*}{*)},
    morekeywords={*,...}
}

\lstset{
    language=Python,
    basicstyle=\ttfamily\footnotesize,
    keywordstyle=\color{blue}\bfseries,
    commentstyle=\color{green!50!black}\itshape,
    stringstyle=\color{red},
    numbers=left,
    numberstyle=\tiny\color{gray},
    stepnumber=1,
    numbersep=5pt,
    backgroundcolor=\color{white},
    showspaces=false,
    showstringspaces=false,
    showtabs=false,
    frame=single,
    tabsize=2,
    captionpos=b,
    breaklines=true,
    breakatwhitespace=false,
    escapeinside={\%*}{*)},
    morekeywords={*,...}
}

% Header and footer configuration
\renewcommand{\sectionmark}[1]{\markboth{#1}{}}
\pagestyle{fancy}
\fancyhf{}
\fancyhead[L]{\textbf{Heap Thrust}}
\fancyhead[C]{\textbf{\nouppercase{\leftmark}}}
\fancyhead[R]{Page \thepage\ of \pageref{LastPage}}
\renewcommand{\headrulewidth}{0.4pt}

% Reduce spacing
\setlength{\columnsep}{0.3in}
\setlength{\parindent}{0pt}
\setlength{\parskip}{0.5em}

% ToC formatting
\setlength{\cftbeforesecskip}{0.2em}
\setlength{\cftbeforesubsecskip}{0.1em}

\begin{document}

% Title page (single column)
\onecolumn
\thispagestyle{empty} % No header on cover page
\begin{center}
    \vspace*{2cm}
    
    {\Huge \textbf{Heap Thrust Reference 2026}}
    
    \vspace{1.5cm}
    
    {\Large ICPC Reference (ESCOM Winter version)}
    
    \vspace{2cm}
    
    {\large \textbf{Team Members:}}
    
    \vspace{0.5cm}
    
    {\large
    Enrique Job Calderón (ksobrenat32)\\
    Luis Daniel Salazar (luisdakan)\\
    Gustavo Valenzuela (GusTimeTraveler)\\
    }
    
    \vspace{2cm}
    
    {\large \textbf{Institution:}}
    
    {\large Facultad de Ingeniería, UNAM - CPCFI}
    
    \vspace{2cm}
    
    {\large \today}
    
\end{center}

\newpage

% Table of Contents
\setcounter{page}{1} % Start page numbering from ToC
\pagestyle{fancy} % Enable header starting from ToC
\begin{multicols}{2}
\tableofcontents
\end{multicols}

\newpage

% Switch to two-column format
\twocolumn

\section{Python stuff}

\subsection{Regular input}

\lstinputlisting[language=Python]{code/Python/input.py}
\subsection{Fast I/O (If we use this, we're down bad)}

\lstinputlisting[language=Python]{code/Python/fastio.py}

\subsection{List comprehensions}
\lstinputlisting[language=Python]{code/Python/list_comprehension.py}

\subsection{Strings}
\lstinputlisting[language=Python]{code/Python/strings.py}


\subsection{Regex}

We must be careful with these functions since they could blow into $O(2^n)$ complexity.

\lstinputlisting[language=Python]{code/Python/regex.py}

\subsection{Itertools}
\lstinputlisting[language=Python]{code/Python/itertools.py}

\subsection{Scripts}

\subsubsection{Combinatorics calculator}

\lstinputlisting[language=Python]{code/Python/combinatorics_calc.py}

\subsubsection{Random array generator}

\lstinputlisting[language=Python]{code/Python/test_arr.py}

\subsubsection{Generate all subsets}
\lstinputlisting[language=Python]{code/Python/gen_subsets.py}

\section{C++ utils}

\subsection{Common STL Algorithms}

\subsection*{Sorting Algorithms}
\begin{tabularx}{\linewidth}{|l|l|X|}
    \hline
    \textbf{Function} & \textbf{Parameters} & \textbf{Description} \\
    \hline
    sort & begin, end, [comp] & Standard unstable sort (O(n log n)) \\
    \hline
    stable\_sort & begin, end, [comp] & Stable sort preserves element order \\
    \hline
    is\_sorted & begin, end, [comp] & Checks if range is sorted (returns bool) \\
    \hline
    nth\_element & begin, nth, end, [comp] & Partitions around nth element \\
    \hline
\end{tabularx}

\subsection*{Searching Functions}
\begin{tabularx}{\linewidth}{|l|l|X|}
    \hline
    \textbf{Function} & \textbf{Parameters} & \textbf{Description} \\
    \hline
    lower\_bound & begin, end, val, [comp] & First element $\leq$ value \\
    \hline
    upper\_bound & begin, end, val, [comp] & First element > value \\
    \hline
    binary\_search & begin, end, val, [comp] & Existence check in sorted range \\
    \hline
    find & begin, end, val & Linear search for value \\
    \hline
    find\_if & begin, end, pred & Find first matching predicate \\
    \hline
\end{tabularx}


\subsection*{Sequence Operations}
\begin{tabularx}{\linewidth}{|l|l|X|}
    \hline
    \textbf{Function} & \textbf{Parameters} & \textbf{Description} \\
    \hline
    reverse & begin, end & Reverse elements in-place \\
    \hline
    rotate & begin, mid, end & Rotate elements left \\
    \hline
    next\_permutation & begin, end & Generate next permutation \\
    \hline
    unique & begin, end, [pred] & Remove consecutive duplicates \\
    \hline
    remove & begin, end, val & Remove elements equal to value \\
    \hline
\end{tabularx}

\subsection*{Numerical Functions}
\begin{tabularx}{\linewidth}{|l|l|X|}
    \hline
    \textbf{Function} & \textbf{Parameters} & \textbf{Description} \\
    \hline
    accumulate & begin, end, init, [op] & Sum/accumulate elements \\
    \hline
    partial\_sum & begin, end, dest, [op] & Compute prefix sums \\
    \hline
    \_\_gcd & a, b & Greatest common divisor (C++17) \\
    \hline
    lcm & a, b & Least common multiple (C++17) \\
    \hline
    iota & begin, end, val & Fill with consecutive values \\
    \hline
\end{tabularx}

\subsection*{Memory/Array Operations}
\begin{tabularx}{\linewidth}{|l|l|X|}
    \hline
    \textbf{Function} & \textbf{Parameters} & \textbf{Description} \\
    \hline
    memset & ptr, value, count & Fill memory with byte value \\
    \hline
    fill & begin, end, value & Fill range with value \\
    \hline
    fill\_n & begin, count, value & Fill N elements with value \\
    \hline
    copy & src\_b, src\_e, dest & Copy range to destination \\
    \hline
    copy\_if & src\_b, src\_e, dest, pred & Copy elements matching predicate \\
    \hline
\end{tabularx}


\subsection*{Utility Functions}
\begin{tabularx}{\linewidth}{|l|l|X|}
    \hline
    \textbf{Function} & \textbf{Parameters} & \textbf{Description} \\
    \hline
    swap & a, b & Swap two values \\
    \hline
    max\_element & begin, end, [comp] & Find maximum element \\
    \hline
    min\_element & begin, end, [comp] & Find minimum element \\
    \hline
    count & begin, end, val & Count element occurrences \\
    \hline
    all\_of & begin, end, pred & Check all elements satisfy condition \\
    \hline
\end{tabularx}

\subsection*{Mathematical / Bitwise Builtins}
\begin{tabularx}{\linewidth}{|l|l|X|}
    \hline
    \textbf{Function} & \textbf{Parameters} & \textbf{Description} \\
    \hline
    \_\_builtin\_popcount & x (int) & Count number of 1-bits \\
    \hline
    \_\_builtin\_popcountll & x (long long) & Count number of 1-bits (64-bit) \\
    \hline
    \_\_builtin\_clz & x (unsigned int) & Count leading zeros \\
    \hline
    \_\_builtin\_clzll & x (unsigned long long) & Leading zeros (64-bit) \\
    \hline
    \_\_builtin\_ctz & x (unsigned int) & Count trailing zeros \\
    \hline
    \_\_builtin\_ctzll & x (unsigned long long) & Trailing zeros (64-bit) \\
    \hline
    \_\_builtin\_parity & x & Return 1 if \#bits is odd \\
    \hline
    \_\_builtin\_ffs & x & Position of least significant 1-bit (1-indexed) \\
    \hline
    \_\_lg & x & Floor of $\log_2(x)$ (index of highest bit) \\
    \hline
\end{tabularx}

\subsection*{Priority Queues and Heaps}
\begin{tabularx}{\linewidth}{|l|l|X|}
    \hline
    \textbf{Function} & \textbf{Parameters} & \textbf{Description} \\
    \hline
    priority\_queue & [type], [container], [comp] & Max-heap by default \\
    \hline
    make\_heap & begin, end, [comp] & Turn range into heap \\
    \hline
    push\_heap & begin, end, [comp] & Push element into heap \\
    \hline
    pop\_heap & begin, end, [comp] & Pop max element into end \\
    \hline
    sort\_heap & begin, end, [comp] & Heap sort \\
    \hline
\end{tabularx}

\subsection*{Set / Map Utilities}
\begin{tabularx}{\linewidth}{|l|l|X|}
    \hline
    \textbf{Operation} & \textbf{Usage} & \textbf{Description} \\
    \hline
    s.lower\_bound(x) & set/map & First element $\geq$ x \\
    \hline
    s.upper\_bound(x) & set/map & First element $>$ x \\
    \hline
    s.equal\_range(x) & multiset/map & Pair of lower/upper bound \\
    \hline
    s.erase(it) & iterator & Erase element at iterator \\
    \hline
    s.find(x) & key & Iterator to key or end \\
    \hline
\end{tabularx}

\subsection*{String Functions}
\begin{tabularx}{\linewidth}{|l|l|X|}
    \hline
    \textbf{Function} & \textbf{Parameters} & \textbf{Description} \\
    \hline
    stoi, stol, stoll & string, [pos], [base] & Convert string $\to$ int/long/ll \\
    \hline
    stoul, stoull & string, [pos], [base] & Convert string $\to$ unsigned \\
    \hline
    stod, stof, stold & string & Convert string $\to$ double/float/long double \\
    \hline
    to\_string & value & Convert number $\to$ string \\
    \hline
    substr & pos, len & Substring \\
    \hline
    find & str, pos & Find first occurrence \\
    \hline
    rfind & str, pos & Find last occurrence \\
    \hline
\end{tabularx}

\subsection*{Random Number Utilities}
\begin{tabularx}{\linewidth}{|l|l|X|}
    \hline
    \textbf{Type / Function} & \textbf{Usage} & \textbf{Description} \\
    \hline
    mt19937 rng & chrono::steady\_clock::now() & Fast random generator \\
    \hline
    uniform\_int\_distribution & dist(a,b)(rng) & Random int in [a,b] \\
    \hline
    shuffle & begin, end, rng & Random shuffle \\
    \hline
\end{tabularx}

\subsection*{Other Useful Utilities}
\begin{tabularx}{\linewidth}{|l|l|X|}
    \hline
    \textbf{Function} & \textbf{Parameters} & \textbf{Description} \\
    \hline
    chrono::high\_resolution\_clock & now() & Get precise current time \\
    \hline
    \_\_int128 & value & 128-bit integer (manual I/O needed) \\
    \hline
    bitset $<N>$ & ops: \&, |, \^, <<, >> & Fixed-size bitset manipulation \\
    \hline
    tuple & get$<i>$(t) & Store and access heterogenous data \\
    \hline
    pair & first, second & Store pair of values \\
    \hline
\end{tabularx}

\subsection{Random number generator}

\lstinputlisting[language=C++]{code/CppUtils/RandomNumbers.cpp}

\section{Binary Search in the Answer}

\lstinputlisting[language=C++]{code/BinarySearch/AllBinSearches.cpp}
% ========================================
% SECTION: Graph Theory
% ========================================
\section{Graph Theory}

\subsection{Depth First Search (DFS)}
Example implementation:
\lstinputlisting[caption=DFS Implementation]{code/GraphTheory/dfs.cpp}

\subsection{Breadth First Search (BFS)}
Example implementation:
\lstinputlisting[caption=BFS Implementation]{code/GraphTheory/bfs.cpp}

% Add more graph theory algorithms as needed

% ========================================
% SECTION: Data Structures
% ========================================
\section{Data Structures}

\subsection{Segment Tree}
Example implementation:
\lstinputlisting[caption=Segment Tree]{code/DataStructures/segment_tree.cpp}

\subsection{Fenwick Tree (BIT)}
Example implementation:
\lstinputlisting[caption=Fenwick Tree]{code/DataStructures/fenwick_tree.cpp}

% ========================================
% SECTION: Range Queries
% ========================================
\section{Range Queries}

\subsection{Sparse Table}
Example implementation:
\lstinputlisting[caption=Sparse Table]{code/RangeQueries/sparse_table.cpp}

% ========================================
% SECTION: Mathematics
% ========================================
\section{Mathematics}

\subsection{Number Theory}

\subsubsection{GCD and LCM}
\begin{align*}
\gcd(a, b) &= \gcd(b, a \bmod b) \\
\text{lcm}(a, b) &= \frac{ab}{\gcd(a, b)}
\end{align*}

\lstinputlisting[caption=GCD Implementation]{code/Math/gcd.cpp}

\subsubsection{Extended Euclidean Algorithm}
For $ax + by = \gcd(a, b)$:
\lstinputlisting[caption=Extended GCD]{code/Math/extended_gcd.cpp}

\subsubsection{Modular Arithmetic}
\begin{align*}
a^{-1} \bmod m &\text{ exists iff } \gcd(a, m) = 1 \\
a^{\phi(m)} &\equiv 1 \pmod{m} \text{ (Euler's theorem)} \\
a^{p-1} &\equiv 1 \pmod{p} \text{ (Fermat's little theorem)}
\end{align*}

\lstinputlisting[caption=Modular Exponentiation]{code/Math/mod_pow.cpp}

\subsection{Combinatorics}

\subsubsection{Binomial Coefficients}
\begin{align*}
\binom{n}{k} &= \frac{n!}{k!(n-k)!} \\
\binom{n}{k} &= \binom{n-1}{k-1} + \binom{n-1}{k} \\
\sum_{k=0}^{n} \binom{n}{k} &= 2^n
\end{align*}

\lstinputlisting[caption=Binomial Coefficients]{code/Math/binomial.cpp}

\subsubsection{Catalan Numbers}
\begin{align*}
C_n &= \frac{1}{n+1}\binom{2n}{n} = \frac{(2n)!}{(n+1)!n!} \\
C_n &= \sum_{i=0}^{n-1} C_i C_{n-1-i} \\
C_0 &= 1, \quad C_{n+1} = \frac{2(2n+1)}{n+2}C_n
\end{align*}

\subsection{Probability}

\subsubsection{Basic Formulas}
\begin{align*}
P(A \cup B) &= P(A) + P(B) - P(A \cap B) \\
P(A|B) &= \frac{P(A \cap B)}{P(B)} \\
P(A \cap B) &= P(A|B) \cdot P(B) = P(B|A) \cdot P(A)
\end{align*}

\subsubsection{Expected Value}
\begin{align*}
E[X] &= \sum_{i} x_i \cdot P(X = x_i) \\
E[aX + b] &= aE[X] + b \\
E[X + Y] &= E[X] + E[Y]
\end{align*}

\subsection{Algebra}

\subsubsection{Matrix Operations}
Matrix multiplication complexity: $O(n^3)$
\begin{align*}
(AB)_{ij} &= \sum_{k=1}^{n} A_{ik}B_{kj} \\
(AB)^T &= B^T A^T \\
\det(AB) &= \det(A)\det(B)
\end{align*}

\subsubsection{System of Linear Equations}
Gaussian elimination: $O(n^3)$
\begin{align*}
Ax &= b \\
\text{If } \det(A) &\neq 0 \text{, unique solution exists}
\end{align*}

\subsection{Calculus}

\subsubsection{Derivatives}
\begin{align*}
(x^n)' &= nx^{n-1} \\
(e^x)' &= e^x \\
(\ln x)' &= \frac{1}{x} \\
(\sin x)' &= \cos x \\
(\cos x)' &= -\sin x
\end{align*}

\subsubsection{Integration}
\begin{align*}
\int x^n \, dx &= \frac{x^{n+1}}{n+1} + C \quad (n \neq -1) \\
\int e^x \, dx &= e^x + C \\
\int \frac{1}{x} \, dx &= \ln|x| + C \\
\int \sin x \, dx &= -\cos x + C \\
\int \cos x \, dx &= \sin x + C
\end{align*}

\subsubsection{Series}
\begin{align*}
\sum_{i=1}^{n} i &= \frac{n(n+1)}{2} \\
\sum_{i=1}^{n} i^2 &= \frac{n(n+1)(2n+1)}{6} \\
\sum_{i=0}^{n} x^i &= \frac{1-x^{n+1}}{1-x} \quad (x \neq 1)
\end{align*}

% ========================================
% SECTION: String Algorithms
% ========================================
\section{String Algorithms}

\subsection{KMP (Knuth-Morris-Pratt)}
Example implementation:
\lstinputlisting[caption=KMP Algorithm]{code/Strings/kmp.cpp}

\subsection{Z-Algorithm}
Example implementation:
\lstinputlisting[caption=Z-Algorithm]{code/Strings/z_algorithm.cpp}

% ========================================
% SECTION: Geometry
% ========================================
\section{Computational Geometry}

\subsection{Basic Formulas}

\subsubsection{Distance and Dot Product}
\begin{align*}
d(P_1, P_2) &= \sqrt{(x_2-x_1)^2 + (y_2-y_1)^2} \\
\vec{a} \cdot \vec{b} &= a_x b_x + a_y b_y = |\vec{a}||\vec{b}|\cos\theta
\end{align*}

\subsubsection{Cross Product}
\begin{align*}
\vec{a} \times \vec{b} &= a_x b_y - a_y b_x \\
|\vec{a} \times \vec{b}| &= |\vec{a}||\vec{b}|\sin\theta
\end{align*}

\subsection{Convex Hull}
Example implementation:
\lstinputlisting[caption=Convex Hull (Graham Scan)]{code/Geometry/convex_hull.cpp}

% ========================================
% SECTION: Dynamic Programming
% ========================================
\section{Dynamic Programming}

\subsection{Longest Increasing Subsequence}
Example implementation:
\lstinputlisting[caption=LIS in O(n log n)]{code/DynamicProgramming/lis.cpp}

\subsection{Knapsack}
Example implementation:
\lstinputlisting[caption=0-1 Knapsack]{code/DynamicProgramming/knapsack.cpp}

\end{document}
